% 香港与数字人 ... 取向短文（简体中文译本）
% 与 places.html / places-hong-kong.tex 同文。专有名词保留。无 em dash。
% 用 tectonic 或 XeLaTeX 编译。
\documentclass[11pt,a4paper]{ctexart}

\usepackage[margin=1in]{geometry}
\usepackage{setspace}
\usepackage{enumitem}
\usepackage[hidelinks]{hyperref}

\setCJKmainfont{Songti SC}
\setmainfont{Times New Roman}

\onehalfspacing
\setlength{\parskip}{0.55em}
\setlength{\parindent}{0pt}
\setlist[itemize]{leftmargin=1.4em, itemsep=0.35em, topsep=0.4em}
\Urlmuskip=0mu plus 1mu

\title{香港与数字人}
\author{Jade Zhao · 赵郎溪}
\date{2026年9月 \\ \small 随笔门：\texttt{zhao-langxi} · 进行中 · 非同行评议}

\begin{document}
\maketitle

\noindent\textit{随笔 · 聚焦香港 · 非同行评议}

看起来像人，并不是产品本身。真正悬而未决的问题，仍是日常里的信任：当一个系统看起来、说话也像人时，什么会让人愿意用它、读懂披露，并在真实工作里继续用下去。
这篇笔记把香港看作那一套技术会公开出现在目录与本地演示里的地方之一。它预告研究，不是搬家倡议，也不是产品发布。

\section*{为什么是香港}

香港夹在面向英语的企业工具习惯，与中文内容生产习惯之间。
本地数字人产品页上，粤语、普通话与英语常一起出现。
这座城市的公共人工智能目录，也把数字人框成部门客户服务与查询接口，而不只是娱乐面孔。

智慧政府创新实验室（Smart Government Innovation LAB，数码港）列入的方案包括：中国移动香港的人工智能数字人客服、腾讯云数字人（Tencent Cloud Digital Human）、Momentum Cloud 的表现型虚拟人引擎（Computer And Technologies），以及 01.AI (HK) Limited 的形象 / 声音克隆方案。
另外，香港生产力促进局（HKPC）的「数码不求人」门户列入第一线 DYXnet 的線靈 / Digital Human Platform，用于人工智能主播与实时助理，并支持普通话、粤语与英语。

\begin{itemize}
  \item Smart LAB · AI Digital human customer service（中国移动香港）：
    \url{https://www.smartlab.gov.hk/en/ai_solutions/a-0073}
  \item Smart LAB · Tencent Cloud Digital Human：
    \url{https://www.smartlab.gov.hk/en/ai_solutions/a-0077}
  \item HKPC Digital DIY · DYXnet Digital Human Platform：
    \url{https://ddiy.hkpc.org/en/node/1519}
  \item DYXnet 产品页（香港）：
    \url{https://www.dyxnet.com/hk/digital-human-platform/}
\end{itemize}

诚实读法：公开营销里，大量仍是播报视频、虚拟主播，以及办事大厅问答，再给知识库或大语言模型加一张脸。
这比纯直播偶像更接近有用的业务界面……但仍不等于在信任、披露与拟人化上做深入工作。

\section*{值得一看的视频}

下列标题与网址已对照 YouTube 的 oEmbed 接口核对。它们是取向演示，不是背书。

\begin{itemize}
  \item DYXnet - Digital Human AI 數字人 video \\
    \url{https://www.youtube.com/watch?v=XVA4cLqBN3g}
  \item 数码港香港官方频道 · 相关 AIGC 公开演示（影视 / 创作者平台，不是政府查询形象）：
    【數碼港智慧生活初創培育計劃：創科啟航 · FizzDragon】創意無界：以 AIGC 連結全球創作者 \\
    \url{https://www.youtube.com/watch?v=Kuav2SDlQw4}
  \item 围绕数字人的粤语创作者文化（消费级工具演示，不是企业案例研究）：
    【最逼真1:1 廣東話虛擬人】20秒教識你打造自己一模一樣の虛擬數字人幫你去拍片 \\
    \url{https://www.youtube.com/watch?v=8t9v9TMb-M0}
\end{itemize}

\section*{中国内地，简要}

边境另一侧，内地数字人平台（例如百度的数字员工 / 曦灵栈，以及华为云 MetaStudio）更明显偏向视频、直播与品牌代言人，智能交互多作为附加能力。
香港公共目录的用语更常是「enquiry」「customer service」与多语种自助终端。
这一对比有用。它并不让任何一处成为「日常信任」的最终答案。

\section*{另外两处枢纽，作对照}

\begin{itemize}
  \item \textbf{伦敦 · Synthesia}……企业视频正走向销售、支持与培训用的交互形象。
    \url{https://jobs.ashbyhq.com/synthesia/e934a7ac-b668-46c0-9953-eed4504658b7}
  \item \textbf{旧金山 / 奥克兰 · Soul Machines}……Digital Workforce 被框成客户体验、人力资源及相关面对面工作流。
    \url{https://www.soulmachines.com/digital-workforce}
\end{itemize}

\section*{这篇笔记不主张什么}

\begin{itemize}
  \item 不是建议搬家、宣称实习，或加入任何具名公司。
  \item 不是市场规模数字或厂商偏好百分比。
  \item 不是临床、卫生系统或实验室管线表述。
  \item 不是声称每一个公共服务形象已经赢得信任。
\end{itemize}

\section*{相关}

这张地图所挂的开放问题见 The question：
\url{https://zhao-langxi.github.io/zhao-langxi/question.html}。
已交付的代码证明见 jadexzhao：
\url{https://jadexzhao.github.io/jadexzhao/}。

\vspace{1.2em}
\noindent\small 独立笔记……非印第安纳大学背书研究，亦非同行评议出版物。\\
Jade Zhao · 赵郎溪 · Digital Humans Project Lead · 预计 2027 年 5 月毕业

\end{document}
